% Generated draft: requires mathematical and source review.
\documentclass[11pt]{article}
\usepackage[margin=25mm]{geometry}
\usepackage{fontspec}
\setmainfont{Latin Modern Roman}
\usepackage{amsmath,amssymb,mathtools,unicode-math}
\setmathfont{Latin Modern Math}
\usepackage{enumitem,needspace,xurl,hyperref}
\setlength{\parindent}{0pt}
\setlength{\parskip}{5pt}
\setlength{\emergencystretch}{3em}
\newcommand{\N}{\mathbb{N}}
\newcommand{\Z}{\mathbb{Z}}
\newcommand{\Q}{\mathbb{Q}}
\newcommand{\R}{\mathbb{R}}
\newcommand{\C}{\mathbb{C}}
\newcommand{\F}{\mathbb{F}}
\newcommand{\A}{\mathbb{A}}
\newcommand{\PP}{\mathbb{P}}
\newcommand{\E}{\mathbb{E}}
\newcommand{\eps}{\varepsilon}
\newcommand{\dd}{\,\mathrm{d}}
\begin{document}
\section*{UnsolvedMath ID: 30002799}
\textbf{Source problem code:} catalogue entry
\subsection{Definitions and mathematical statement}
The source poses two questions concerning ordinary and symbolic powers. It does not specify the ground field, the admissible ambient rings for the second question, or the symbolic-power convention for nonradical ideals. The formulations below therefore retain these qualifications explicitly.

For a proper ideal \(I\) of a commutative Noetherian ring \(R\), its ordinary power \(I^m\) is the ideal generated by products of \(m\) elements of \(I\). Under the associated-prime convention, set
\[
 I^{(m)}=\bigcap_{\mathfrak p\in\operatorname{Ass}_R(R/I)}\bigl(I^mR_{\mathfrak p}\cap R\bigr),
 \qquad
 \operatorname{bight}(I)=\max_{\mathfrak p\in\operatorname{Ass}_R(R/I)}\operatorname{ht}_R(\mathfrak p),
 \quad m\in\Z_{\geq1},
\]
where contraction to \(R\) is understood in the displayed intersections. A minimal-prime convention instead uses \(\operatorname{Min}_R(R/I)\); these choices agree for radical ideals in this setting but need not agree for general ideals. The source does not resolve which convention or ideal class is intended in its second question.

For the first question, the usual finite reduced point-set interpretation uses an unspecified field \(k\), the bihomogeneous coordinate ring
\[
 S=k[x_0,x_1,y_0,y_1],\qquad
 \deg(x_i)=(1,0),\quad \deg(y_i)=(0,1),
\]
and the bihomogeneous vanishing ideal \(I(X)\subseteq S\) of a finite reduced set \(X\subseteq\PP^1_k\times\PP^1_k\). Here arithmetically Cohen--Macaulay means that \(S/I(X)\) is Cohen--Macaulay; equivalently, its depth at the homogeneous maximal ideal equals its Krull dimension. This interpretation, including finiteness and reducedness, requires confirmation from the original report: the supplied statement says only ``Cohen--Macaulay point sets.''

\emph{Geometric characterization question.} Under this interpretation, characterize exactly those arithmetically Cohen--Macaulay point sets \(X\) for which
\[
 I(X)^3=I(X)^{(3)}.
\]
The source also describes this as a question about equivalent line unions in \(\PP^3\). For a set of \(k\)-rational points, one possible precise correspondence is obtained by giving all four variables of \(S\) degree one: the ideal of a point \((a,b)\) is generated by one linear form in the \(x\)-variables and one in the \(y\)-variables, and hence defines a line in \(\operatorname{Proj}(S)=\PP^3_k\). The same intersection of point ideals then defines the corresponding union of lines. The source does not specify this correspondence, rationality assumptions, or the intended class of line unions; its equivalence must be understood subject to that unresolved qualification.

\emph{Delayed failure question.} Does the intended class of ideals contain, for every \(t\in\Z_{\geq1}\), an ideal \(I_t\subseteq R_t\), with \(h_t=\operatorname{bight}(I_t)\), such that
\[
 \bigl(\forall m\in\Z_{\geq1},\ m\leq h_t+t-1
       \Longrightarrow I_t^m=I_t^{(m)}\bigr)
 \qquad\text{and}\qquad
 I_t^{h_t+t}\ne I_t^{(h_t+t)}?
\]
The admissible rings \(R_t\), whether they may vary with \(t\), and any requirements that the ideals be radical, homogeneous, or geometrically defined are not supplied. This is consequently a quantified question schema pending specification of that class, rather than a fully determined existence assertion.
\subsection{Short English statement}
Characterize the Cohen--Macaulay point sets in \(\PP^1\times\PP^1\) whose third ordinary and symbolic ideal powers coincide. Also ask whether, for every positive integer \(t\), an ideal can have equality through exponent \(\operatorname{bight}(I)+t-1\) and first failure at the next exponent. The field, point-set interpretation, line-union correspondence, and admissible ideal class require clarification. The supplied open status and literature assessment are dated curation, not independently verified conclusions.
\subsection{Sources}
\begin{itemize}
\item Local UnsolvedMath registry: supplied complete record for permanent ID \(30002799\), source problem code \texttt{catalogue entry}, title ``Ordinary and Symbolic Powers of Point-Set Ideals.'' The record belongs to the Oberwolfach Reports Open Problems collection and gives proposed year 2015; no individual proposer is supplied.
\item \emph{Mini-Workshop: Ideals of Linear Subspaces, Their Symbolic Powers and Waring Problems} (2015), Oberwolfach Reports, pp.~489--532, DOI \texttt{10.4171/owr/2015/9}. Supplied primary report reference: \url{https://doi.org/10.4171/owr/2015/9}. No exact problem page or speaker attribution is supplied.
\item Supplied literature reference: \emph{Points fattening on P1 times P1}, \url{https://doi.org/10.1016/j.jpaa.2013.12.007}. The record supplies no authors or theorem numbers.
\item Literature assessment checked at \texttt{2026-08-21}: the supplied curation labels the problem open and reports that, for ACM point sets in \(\PP^1\times\PP^1\), equality of symbolic and ordinary powers is characterized through the cubic case, while the requested general geometric characterization and ideal-family threshold remain unsettled. This is supplied dated curation derived from the OWR material, not an independently verified literature assessment. The background expressly requests independent verification of time-sensitive status and cited claims.
\item The registry labels the statement \texttt{corrected\_verified}, explaining that mathematical content was preserved while notation, typography, or omitted definitions were repaired. This is a supplied editorial label, not independent mathematical verification of this draft.
\end{itemize}
\end{document}
